% Part: computability
% Chapter: computability-theory
% Section: ce-sets

\documentclass[../../../include/open-logic-section]{subfiles}

\begin{document}

\olfileid{cmp}{thy}{ces}
\olsection{గణనీయంగా లెక్కించదగిన సమితులు}

\begin{defn}
ఒక సమితి శూన్య సమితి అయినా, లేదా ఒక గణనీయ ప్రమేయపు వ్యాప్తి అయినా,
దాన్ని \emph{గణనీయంగా లెక్కించదగిన సమితి} అంటాం.
\end{defn}

\begin{history}
గణనీయంగా లెక్కించదగిన సమితులను బదులుగా \emph{పునరావృత్తంగా
లెక్కించదగిన సమితులు} అని కూడా అంటారు. ఇదే అసలు పరిభాష. నేడు రెండు
పేర్లూ, అలాగే ``c.e.'', ``r.e.'' అనే సంక్షిప్త రూపాలూ సాధారణంగా
వాడుకలో ఉన్నాయి.
\end{history}

\begin{explain}
ఈ నిర్వచనం ఏమి చెబుతుందో, ఈ పరిభాష ఎందుకు తగినదో ఆలోచించాలి. $S$
అనేది గణనీయ ప్రమేయం~$f$ యొక్క వ్యాప్తి అయితే,
\[
S=\{f(0),f(1),f(2),\dots\},
\]
కాబట్టి $f$ అనేది $S$లోని మూలకాలను ``లెక్కిస్తుంది'' అని భావించవచ్చు.
నిర్వచనం ప్రకారం $f$ వృద్ధి ప్రమేయం కానవసరం లేదని, అంటే లెక్కింపు
వృద్ధి క్రమంలో ఉండనవసరం లేదని గమనించండి. నిజానికి $f$ ఒకటి-ఒకటి
ప్రమేయం కూడా కానవసరం లేదు; అంటే $S$ యొక్క $f(0)$, $f(1)$,
$f(2)$, \dots{} లెక్కింపులో పునరావృత్తులను అనుమతిస్తాం. ఉదాహరణకు,
స్థిర ప్రమేయం $f(x)=0$ అనేది $\{0\}$ సమితిని లెక్కిస్తుంది.
\end{explain}

ప్రతి గణనీయ సమితీ గణనీయంగా లెక్కించదగినదే. దీన్ని చూడటానికి, $S$
గణనీయం అనుకుందాం. $S$ శూన్య సమితి అయితే, నిర్వచనం ప్రకారమే అది
గణనీయంగా లెక్కించదగినది. లేకపోతే, $S$లో ఏదైనా మూలకం $a$ను తీసుకోండి.
$f$ను ఇలా నిర్వచించండి:
\[
f(x)=
\begin{cases}
x & \text{$\Char{S}(x)=1$ అయితే} \\
a & \text{ఇతర సందర్భాల్లో.}
\end{cases}
\]
అప్పుడు $f$ ఒక గణనీయ ప్రమేయం; $S$ అనేది~$f$ వ్యాప్తి.

\end{document}
